\section{The node and two nonreduced deformation spaces}\label{sec:examples}

\subsection*{The period at a node}
Write the deformation of the ordinary double point as
\[
 x_1y_1=z,\qquad x_2y_2=z-s.
\]
On the open torus the relative volume form is
$\Omega=dz\wedge d\log x_1\wedge d\log x_2$.
For $s>0$ the equations $y_1=\overline{x}_1$ and
$y_2=-\overline{x}_2$ cut out the vanishing sphere
$|x_1|^2+|x_2|^2=s$. The map to $z=|x_1|^2$ fibres the interior
of the interval $[0,s]$ by two-tori, with a different circle collapsing
at either end. Orient the sphere by $dz\wedge d\arg x_1\wedge d\arg x_2$.
For a holomorphic unit $g=g(z,s)$ the two angular integrations give
\begin{equation}\label{ex:nodeintegral}
 \int_{S^3}g\Omega=(2\pi i)^2\int_0^s g(z,s)\,dz.
\end{equation}
The Milnor fibre is a cotangent disc bundle of $S^3$, whose boundary
is $S^3\times S^2$ since $S^3$ is parallelizable. Inclusion of the
boundary induces an isomorphism on $H_3$. Thus
\eqref{ex:nodeintegral} evaluates the marked link class on a generator.
For $g=1+cz$ this period is $(2\pi i)^2(s+cs^2/2)$. Even on a smooth
local base, the choice of relative form can give a nonlinear period
coordinate.

\subsection*{A length-ten base}
Consider the hexagon with cyclic primitive edge vectors
\[
 (1,0),\ (1,1),\ (0,1),\ (-4,1),\ (1,-2),\ (1,-1).
\]
For edge vectors $d_i$ write $M_k=\sum_i d_i t_i^k$. For an $n$-gon
the moments through degree $n-1$ generate the full ideal. Indeed,
$M_0=0$ and the polynomial $\prod_i(T-t_i)$ gives a recurrence of
order $n$ for $M_k$. Here $t_6=0$ and $M_1=0$ give
\[
 (t_1,\ldots,t_6)=(x+5y-3z,-x-y+2z,x,y,z,0).
\]
Let $I\subset\C[x,y,z]$ be the resulting moment ideal. In graded
lexicographic order $x>y>z$, a Gr\"obner basis is
\[
\begin{gathered}
 5x^2-7xz-5y^2+5yz-z^2,\qquad
 5xy-3xz+10y^2-15yz+6z^2,\\
 3xz^2+5yz^2-z^3,\qquad y^3+yz^2-z^3,\qquad 2y^2z-z^3,\\
 yz^3,\qquad z^4.
\end{gathered}
\]
It is obtained from the two components of $M_2$ and $M_3$.
Every quartic monomial is divisible by one of its leading monomials.
By homogeneity, all moments of higher degree vanish in this quotient.
The standard monomials in degrees
two and three are $xz,y^2,yz,z^2$ and $yz^2,z^3$, respectively.
The base is supported at the origin and its Hilbert function is $(1,3,4,2)$.

On this base one can see the restriction imposed by a common change of
height. The shear $x\mapsto x+z^2$ fixing $y,z$ induces the identity on
the tangent space. For $q=x^2+xy-2xz+y^2-2yz+z^2\in I$, however,
the normal form of its pullback is
\[
 q(x+z^2,y,z)\equiv -(7yz^2+4z^3)/3\pmod I.
\]
Both monomials on the right are standard. The shear therefore fails to
restrict to the base. In the same coordinates the substitution
\[
 t_i\longmapsto H(t_i),\qquad
 H(T)=(1+x-y)T+(\alpha+z)T^2+\beta T^3,\quad \alpha,\beta\in\C,
\]
defines an automorphism. Indeed each $H(t_i)^k$ has an expansion in
$t_i$ with coefficients independent of $i$, and its moment vanishes.
The linear part of the induced endomorphism is the identity, so it is
invertible on the finite local algebra. The two edge rows give, for
the divided differences $U_i=(H(\zeta)-H(t_i))/(\zeta-t_i)$,
\[
 U_1U_2U_5U_6=U_4^4,\qquad U_2U_3U_4=U_5^2U_6.
\]
These are the identities that lift the substitution to the family.

\subsection*{The nilradical of the octagon}
For the octagon of \cite[\S9.3]{Alt} the two moment equations equate
the power sums of the triples
\[
 (t_1,t_2,t_8),\ (t_4,t_5,t_6)\qquad\text{and}\qquad
 (t_2,t_3,t_4),\ (t_6,t_7,t_8),
\]
respectively. Degrees one through three suffice by Newton's identities.
Set
\[
 (t_1,\ldots,t_8)=(a+2b+c-e,-a-b+d+e,a,b,c,d,e,0).
\]
The two linear equations now hold. Write $I\subset\C[a,b,c,d,e]$
for the four remaining equations, in degrees two and three.
The three components of the reduced base have ideals
\[
 J_1=(a-e,b),\qquad J_2=(a-d,b,c-e),\qquad J_3=(a,b-e,c).
\]
Their dimensions are $3,2,2$. To compute the thickening along their
union, put
\[
 u=ab,\qquad v=b(b-e),\qquad w=bc,\qquad
 f=(a-e)(a-d),\quad h=(a-e)(c-e).
\]
Then $J_1\cap J_2=(b,f,h)$. Modulo $J_3$, the elements $b,f,h$
become $e,ed,e^2$. It follows that
\[
 J_1\cap J_2\cap J_3=(f-db,h-eb,u,v,w).
\]
Indeed $f-db$ and $h-eb$ belong to $J_3$ and $b$ is a
nonzerodivisor modulo $J_3$. The two quadratic moment equations give
$f-db=-u-v$ and $h-eb=-2u-v-2w$ modulo $I$.

It remains to determine the powers of $(u,v,w)$ in the quotient.
Division by a Gr\"obner basis of the four moment equations, in graded
lexicographic order $a>b>c>d>e$, gives
\[
 u(u+v)=v(u+v)=uw=vw=v^3=w^3=0,\qquad
 \operatorname{NF}(v^2)=b^2d^2-b^2de-bd^2e+bde^2.
\]
Here $\operatorname{NF}$ denotes the polynomial normal form for that
order. The product relations kill every cubic in $u,v,w$. The displayed
quartic is nonzero. Since $I$ is contained in each
$J_i$, this proves that $\sqrt I/I=(u,v,w)$ and that its nilpotency
index is three. Both polygon calculations concern the local deformation
space. In a compact threefold its period image is further intersected
with the subspace of global homological relations.
